\begin{pastexam}[label={px:ntu-108-esoe-ds-05}]{臺大工科海洋所 108 學年度 第 5 題}
（程式設計；本卷規定程式題用 C++ 或 Java）

\textbf{(15 \%)} Write a recursive method \textsf{preorderElementAt(int k)} that returns the $k$th
element in the preorder listing of the subtree in a binary tree. For example, if $k = 1$, return the
element in the root node. The method should run in O($d$) time, where $d$ is the depth of the
subtree. If the subtree doesn't have $k$ nodes, return null. Define all necessary classes and
associated fields to accomplish the function.
\end{pastexam}

\begin{pastsoln}
關鍵是「$O(d)$」：如果只有左右指標，要找前序第 $k$ 個只能真的走訪，最壞要看 $k$ 個節點，不是 $O(d)$。
所以要替每個節點多存一個欄位 $\mathit{size}$，記以它為根的子樹有幾個節點（資結書第~9.4~節做名次查詢
用的也是這個欄位）。 題目說的 depth of the subtree 指這棵子樹的高度（最深的節點有多深）。

\textbf{想法}。 前序走訪的順序是：根，然後整棵左子樹，然後整棵右子樹。 令左子樹有 $L$ 個節點：
\begin{itemize}
\item $k = 1$：答案是根。
\item $2 \le k \le L + 1$：答案在左子樹，是左子樹前序的第 $k - 1$ 個。
\item $k > L + 1$：答案在右子樹，是右子樹前序的第 $k - 1 - L$ 個。
\end{itemize}
每次只往一個子樹走一層，不必看另一邊。

卷首規定程式題用 C++ 或 Java，下面用 C++。 Java 的 \texttt{null} 在這裡用 \texttt{nullptr} 表示：
函式回傳元素的指標，找不到時回傳 \texttt{nullptr}。 元素型別用 \texttt{int}，換成樣板也一樣。

\begin{minted}[fontsize=\small]{cpp}
class BinaryTreeNode {
public:
    int item;
    BinaryTreeNode* left;
    BinaryTreeNode* right;
    int size;                  // 以這個節點為根的子樹有幾個節點

    BinaryTreeNode(int x, BinaryTreeNode* l = nullptr, BinaryTreeNode* r = nullptr)
        : item(x), left(l), right(r) {
        size = 1 + sizeOf(l) + sizeOf(r);
    }

    static int sizeOf(const BinaryTreeNode* t) {
        return t == nullptr ? 0 : t->size;
    }

    // 回傳這棵子樹前序走訪的第 k 個元素；沒有第 k 個時回傳 nullptr
    int* preorderElementAt(int k) {
        if (k < 1 || k > size) {
            return nullptr;
        }
        if (k == 1) {
            return &item;                      // 前序的第一個是根
        }
        int leftSize = sizeOf(left);
        if (k - 1 <= leftSize) {
            return left->preorderElementAt(k - 1);
        }
        return right->preorderElementAt(k - 1 - leftSize);
    }
};
\end{minted}

整棵樹的類別只要存根，查詢時呼叫 \texttt{root->preorderElementAt(k)}（樹是空的時直接回傳
\texttt{nullptr}）。 樹在插入、刪除節點時要維護 $\mathit{size}$：從被改的位置往根走，路上每個節點的
$\mathit{size}$ 加一或減一，這一段也是 $O(d)$，不會讓那些操作變慢。

\textbf{正確性}。 對子樹的高度做歸納。 第一個 \texttt{if} 處理子樹不到 $k$ 個節點（以及 $k < 1$）的情形，
回傳 \texttt{nullptr}，符合題目。 之後一定有 $1 \le k \le \mathit{size}$：
\begin{itemize}
\item $k = 1$ 回傳根，對。
\item $2 \le k \le L + 1$：這時 $L \ge k - 1 \ge 1$，左子樹不是空的，不會讀到空指標；前序列表的第
$2$ 到第 $L + 1$ 個恰好是左子樹的前序列表，所以第 $k$ 個是左子樹的第 $k - 1$ 個。 由歸納假設，
遞迴呼叫回傳的就是它。
\item $k \ge L + 2$：$k - 1 - L \ge 1$，而且 $k - 1 - L \le \mathit{size} - 1 - L$ 等於右子樹的節點數，所以右子樹
不是空的，而且有第 $k - 1 - L$ 個；前序列表從第 $L + 2$ 個開始是右子樹的前序列表，所以答案是右子樹的
第 $k - 1 - L$ 個。
\end{itemize}

\textbf{複雜度}。 每次呼叫只做常數的工作，再往下一層呼叫一次，所以呼叫次數不超過子樹的高度加一，
時間 $O(d)$；遞迴深度也是 $O(d)$。 空間上每個節點多一個整數。

例：根 $A$，左子樹是 $B$（左 $D$、右 $E$），右子樹是 $C$（右 $F$）；前序是 $A, B, D, E, C, F$。
找 $k = 5$：根的 $L = 3$，$5 > 4$，到右子樹找第 $5 - 1 - 3 = 1$ 個，就是 $C$。 找 $k = 4$：$4 \le 4$，到
左子樹找第 $3$ 個；$B$ 的 $L = 1$，$3 > 2$，到 $B$ 的右子樹找第 $3 - 1 - 1 = 1$ 個，就是 $E$。 找 $k = 7$：
$7 > \mathit{size} = 6$，回傳 \texttt{nullptr}。
\end{pastsoln}
